\begin{theorem}[Triangle Inequality]
Let \((V,\langle\cdot,\cdot\rangle)\) be an inner‑product space and let \(x,y\in V\). Then
\[
\|x+y\|\le \|x\|+\|y\|,
\]
where the norm is defined by \(\|v\|=\sqrt{\langle v,v\rangle}\) for every \(v\in V\).
\end{theorem}

\begin{proof}
\textbf{Norm definition.}  
For any \(v\in V\),
\[
\|v\|:=\sqrt{\langle v,v\rangle}\ge 0,
\]
since \(\langle v,v\rangle\) is a non‑negative real number.

\medskip

\textbf{Zero‑vector cases.}  
If \(x=0\) then \(\|x+y\|=\|y\|=\|x\|+\|y\|\); similarly if \(y=0\).  
Henceforth assume \(x\neq0\) and \(y\neq0\).

\medskip

\textbf{Expansion of \(\|x+y\|^{2}\).}  
Using linearity in the first argument and conjugate symmetry,
\begin{align*}
\langle x+y,x+y\rangle
&=\langle x,x\rangle+\langle x,y\rangle+\langle y,x\rangle+\langle y,y\rangle\\
&=\|x\|^{2}+\|y\|^{2}+ \bigl(\langle x,y\rangle+\overline{\langle x,y\rangle}\bigr)\\
&=\|x\|^{2}+2\operatorname{Re}\langle x,y\rangle+\|y\|^{2}.
\end{align*}
Consequently,
\begin{equation}
\|x+y\|^{2}= \|x\|^{2}+2\operatorname{Re}\langle x,y\rangle+\|y\|^{2}.
\tag{1}
\end{equation}

\medskip

\textbf{Cauchy–Schwarz inequality.}  
For any vectors \(x,y\) in an inner‑product space,
\begin{equation}
|\langle x,y\rangle|\le \|x\|\,\|y\|,
\tag{2}
\end{equation}
with equality iff \(x\) and \(y\) are linearly dependent.  
Since \(\operatorname{Re}z\le|z|\) for any complex number \(z\), we obtain
\begin{equation}
\operatorname{Re}\langle x,y\rangle\le |\langle x,y\rangle|\le \|x\|\,\|y\|.
\tag{3}
\end{equation}

\medskip

\textbf{Upper bound for \(\|x+y\|^{2}\).}  
Substituting \eqref{3} into \eqref{1} yields
\begin{align}
\|x+y\|^{2}
&\le \|x\|^{2}+2\|x\|\,\|y\|+\|y\|^{2}\\
&=(\|x\|+\|y\|)^{2}.
\tag{4}
\end{align}

\medskip

\textbf{Taking square roots.}  
Both sides of \eqref{4} are non‑negative, and the square‑root function is monotone on \([0,\infty)\); therefore
\[
\|x+y\|\le \|x\|+\|y\|.
\]

\medskip

\textbf{Equality condition.}  
Equality in the triangle inequality occurs precisely when equality holds in \eqref{3}.  
Equality in Cauchy–Schwarz requires \(y=\lambda x\) for some \(\lambda\in\mathbb C\); equality in \(\operatorname{Re}z\le|z|\) forces \(\langle x,y\rangle\) to be a non‑negative real number, i.e. \(\lambda\ge0\).  
Thus equality holds iff either one of the vectors is zero or there exists \(\lambda\ge0\) such that \(y=\lambda x\). \(\qed\)
\end{proof}